\documentclass[10pt,a4paper]{amsart}
\usepackage[margin=19mm]{geometry}
\usepackage{amsmath,amssymb,mathtools}
\usepackage[scr=rsfs,cal=euler]{mathalfa}
\usepackage{xfrac}
\usepackage[unicode,hidelinks,bookmarks=false]{hyperref}
\newcommand{\ie}{i.e.\ }
\newcommand{\cf}{cf.\ }
\newcommand{\resp}{respectively\ }
\newcommand{\Subsection}{\subsection}
\providecommand{\nobreakdash}{\nobreak\hbox{-}}
\setlength{\emergencystretch}{2em}
\multlinegap0pt
\usepackage{bm}
\usepackage{bbm}
\usepackage{dsfont}
\usepackage{xspace}
\usepackage{cancel}
\usepackage{mathtools}
\usepackage{cleveref}
\usepackage{amsthm}
\makeatletter
\@ifundefined{remark}{\newtheorem{remark}{Remark}}{}
\@ifundefined{theorem}{\newtheorem{theorem}{Theorem}}{}
\@ifundefined{lemma}{\newtheorem{lemma}[theorem]{Lemma}}{}
\makeatother
\newcommand{\FF}{\ensuremath{\mathcal{F}}}
\newcommand{\FS}[2][]{\ensuremath{\mathrm{FS}^{#1}\left(#2\right)}}
\newcommand{\N}{\ensuremath{\mathbb{N}}}
\newcommand{\PP}{\ensuremath{\mathcal{P}}}
\newcommand{\Z}{\ensuremath{\mathbb{Z}}}
\newcommand{\abs}[1]{\left | #1 \right |}
\newcommand{\ind}{\ensuremath{ \mathbbm{1}}}
\title[On the partition regularity of arithmetic progressions and l]{On the partition regularity of arithmetic progressions and linear equations in k-IP-sets}
\author{Rapha\"el Giordano}
\date{Source: arXiv:2509.17013v1 (CC BY 4.0)}
\begin{document}
\maketitle

\noindent\emph{Source and licence.} On the partition regularity of arithmetic progressions and linear equations in k-IP-sets. Version arXiv:2509.17013v1 (CC BY 4.0), distributed under the Creative Commons Attribution 4.0 International licence, as recorded in the arXiv licence metadata for this exact version and shown on the abstract page https://arxiv.org/abs/2509.17013v1. Full rights notice: licenses/arxiv-2509.17013-CC-BY-4.0.txt.\par\smallskip

\noindent\textit{Editorial scope.} Counted unit: the Remark whose source label is \texttt{sparse\_enough} in the article (source0.tex lines 416--418, source label \texttt{sparse\_enough}), delivered together with the article's own complete original proof of it (source0.tex lines 420--456, one \texttt{proof} environment of 37 lines). No mathematics is written, repaired or re-derived: statement and proof are the article's own text line for line. Required context retained uncounted: neat\_fact (lines 374--381); classification\_regular\_systems\_IP (lines 390--397); uniqueness B negative (lines 403--405). Every definition, display and label of those lines is retained unchanged. Omitted from this package and not redistributed: 1--373, 382--389, 398--402, 406--415, 419--419, 457--704. Nothing inside a retained excerpt is omitted or altered: every retained body is delivered exactly as the article's lines are written, so no omission receipt of any kind accompanies this package. Citations: the retained excerpts carry no citation command at all, so no bibliography entry of the article is transcribed and no reference is redistributed. Reference adaptation: none was needed and none was made; every reference inside the delivered unit resolves inside it, so no unresolved reference or citation is produced. Printed slips kept literal: no printed word, symbol, hypothesis or number of the counted statement or of its proof was altered. Layout and class: the article's own class is \texttt{article} with packages including geometry, amssymb, amsmath, amsthm, babel, url, stmaryrd, bbm, subcaption, mathtools, enumitem, hyperref, cleveref; the wrapper is the lane's pinned \texttt{amsart} editorial wrapper at 10pt on A4, which restates the theorem-like environments the retained text needs and adds the article's own macro definitions that the retained text needs (\texttt{\textbackslash FF}, \texttt{\textbackslash FS}, \texttt{\textbackslash N}, \texttt{\textbackslash PP}, \texttt{\textbackslash Z}, \texttt{\textbackslash abs}, \texttt{\textbackslash ind}), with extra wrapper packages (bm, bbm, dsfont, xspace, cancel, mathtools, cleveref). The delivered numbering is the unit's own. Assets: no retained span contains a figure environment, a graphics-inclusion command, a PDF-page inclusion, a table or a TikZ picture; no image file is copied into this package, the package declares no asset dependency and no image file is redistributed with it. Licence: version arXiv:2509.17013v1 of this article is distributed under the Creative Commons Attribution 4.0 International licence, as recorded in the arXiv licence metadata for this exact version and shown on the abstract page https://arxiv.org/abs/2509.17013v1. A full rights notice is delivered at licenses/arxiv-2509.17013-CC-BY-4.0.txt. Characters: every retained excerpt is pure ASCII, so the delivered engine, XeLaTeX with its default Latin Modern text font, reports no missing character.
\begin{lemma}\label{neat_fact}
    Let $\FF \subseteq \PP(\N)$. The following statements are equivalent:

    \begin{enumerate}
        \item $\FF$ is partition regular on every $IP$-set, i.e. for every finite coloring of an IP-set there exists a monochromatic subset which belongs to $\FF$.
        \item Every $IP$-set contains a subset $A \in \FF$.
    \end{enumerate}
\end{lemma}

\begin{theorem}[Rado's theorem on IP-sets]\label{classification_regular_systems_IP}
    Let $A \in \Z^{l\times n}$ with columns $\Vec{a_1}, \ldots, \Vec{a_n} \in \Z^l$. The following statements are equivalent:
    \begin{enumerate}
        \item The system of equations in matrix form $A\Vec{x} = \Vec{0}$, where $\Vec{x}= (x_1,\ldots,x_n)$, is partition regular on every IP-set.
        \item There exist $m\in \N$ and $I_1, \ldots, I_m \subseteq \{1,\ldots,n\}$ such that 
        $$\bigcup_{j=1}^m I_j = \{1,\ldots,n\} \text{ and } \forall j \in \{1,\ldots,m\} \; \sum_{i \in I_j} \Vec{a_i} = \Vec{0}.$$
    \end{enumerate}
\end{theorem}

        \begin{lemma}\label{uniqueness B negative}
            If $\sum_{i=1}^{m} \gamma_i B^i = \sum_{i=1}^{m} \delta_i B^i$ for some $B\in \N$ and $\gamma_i, \delta_i \in \Z$ such that $\abs{\gamma_i},\abs{\delta_i} < \frac{B}{2}$, then for every $i \in \{1,\ldots,m\}$, $\gamma_i = \delta_i$.
        \end{lemma}

    \begin{remark}\label{sparse_enough}
        In \cref{uniqueness B negative}, we can replace the powers $B^i$ by a sequence of numbers $(y_i)_{i=1}^m$ which increases sufficiently fast.
    \end{remark}

\begin{proof}[Proof of \cref{classification_regular_systems_IP}]
    $(2) \implies (1):$ Let $I_1,\ldots, I_m \subseteq \{1,\ldots,n\}$ such that 
        $$\bigcup_{j=1}^m I_j = \{1,\ldots,n\} \text{ and } \forall j \in \{1,\ldots,m\} \; \sum_{i \in I_j} \Vec{a_i} = \Vec{0}.$$

        In light of \cref{neat_fact}, we show that $\FS{(y_i)_{i=1}^m}$ contains a solution for any $y_1,\ldots,y_m \in \N$. We set 
        $$x_i = \sum_{j=1}^m \ind_{I_j}(i) y_j.$$
        Notice that $x_i \neq 0$ because $i\in I_j$ for some $j$ by assumption. Then we have
    \begin{align*}
        A\Vec{x} &= \sum_{i=1}^n x_i \Vec{a_i}
            = \sum_{i=1}^n \left (\sum_{j=1}^m \ind_{I_j}(i) y_j \right) \Vec{a_i}
            = \sum_{j=1}^m y_j \left(\sum_{i=1}^n \ind_{I_j}(i) \Vec{a_i}\right)\\
            &= \sum_{j=1}^m y_j\underbrace{\left (\sum_{i\in I_j} \Vec{a_i} \right )}_{=\Vec{0}}
            =\Vec{0}
    \end{align*}

    
    $(1) \implies (2):$ By partition regularity, there exist $m \in \N, x_1, \ldots, x_n \in \FS{(B^j)_{j=1}^m}$ where
    $$B= 2n  \max\left\{\abs{a_{i,j}}: 1 \leq i\leq n, \; 1 \leq j \leq m \right\} + 1$$
    such that
    $$A\Vec{x}= x_1 \Vec{a_1} + \cdots + x_n \Vec{a_n} =\Vec{0}.$$
    Write $$x_i= \sum_{j=1}^m \epsilon_{i,j} B^j$$ with $\epsilon_{i,j} \in \{0,1\}$ for $i=1,\ldots,n,\; j= 1,\ldots,m$. Then
    $$\sum_{i=1}^n \left( \sum_{j=1}^m \epsilon_{i,j} B^j\right) \Vec{a_i} = \Vec{0}$$
    and thus
    $$\sum_{j=1}^m B^j\left (\sum_{i=1}^n \epsilon_{i,j} \Vec{a_i} \right )= \Vec{0}.$$
    
    In particular for each line $ t \in \{1,\ldots,l\}$, we get
    $$\sum_{j=1}^m B^j\left (\sum_{i=1}^n \epsilon_{i,j} a_{i,t} \right )= 0.$$
    Since $$\abs{\sum_{i=1}^n   \epsilon_{i,j} a_{i,t}} \leq n \max\{\abs{a_{i,j}}: 1 \leq i\leq n, \; 1 \leq j \leq m\} < \frac{B}{2},$$
    it follows by \cref{uniqueness B negative} that for every $j\in \{1,\ldots,m\}$,
    $$\sum_{i=1}^n \epsilon_{i,j} a_{i,t} = 0.$$
    The equalities on each line imply
    $$\sum_{i=1}^n \epsilon_{i,j} \Vec{a_i}= \Vec{0}.$$
    
    For $j=1,\ldots,m$, let $$I_j = \{i \in \{1,\ldots,n\}: \; \epsilon_{i,j}= 1\}.$$ By construction,
    $$\sum_{i\in I_j} \Vec{a_i} = \sum_{i=1}^n  \epsilon_{i,j} \Vec{a_i} = \Vec{0}.$$
    Let $i \in \{1,\ldots,n\}$. Since $x_i \neq 0$, there exists $j\in \{1,\ldots,m\}$ such that $\epsilon_{i.j}= 1$ which implies that $i \in I_j$. Therefore $$\bigcup_{j=1}^m I_j = \{1,\ldots,n\}$$ and the criterion is satisfied.    
\end{proof}

\end{document}
